\umaappendix{Prompts del sistema}
\label{ap:prompts}

Este apéndice recoge el texto literal de los \emph{prompts} de las \emph{pipelines} y de la ingesta; el resto (contexto del usuario, Arena y análisis del historial) está en el mismo fichero del repositorio \cite{repositorio}. Todos están definidos en un único fichero, \path{src/AsistenteAyuntamiento.Application/Common/Prompts/PromptCatalog.cs}, de modo que ningún servicio escribe instrucciones propias. Los \emph{prompts} evolucionan con el sistema; los textos de este apéndice se extraen de ese fichero automáticamente con un script (\path{memoria/scripts/extraer-prompts.py}), por lo que reproducen exactamente el código.

Los marcadores \texttt{\{0\}}, \texttt{\{1\}} y \texttt{\{2\}} se sustituyen en tiempo de ejecución por los valores que se indican en cada caso.

\lstset{numbers=none,basicstyle=\ttfamily\footnotesize,breakindent=0pt}

\section{Pipeline básica}

La pipeline básica, y con ella las variantes recursiva y semántica, genera las respuestas con el modelo configurado por cada ayuntamiento. Elige el \emph{prompt} de sistema según la búsqueda haya encontrado documentos o no.

\subsection{Respuesta en \emph{streaming} con documentos}

Es el caso habitual de la pipeline básica. El \emph{prompt} de sistema va seguido, si el usuario los tiene, de sus temas y zonas de interés, con una indicación final (\texttt{UserContextGuidance}) que pide usarlos para priorizar o enfocar la respuesta.

\lstinputlisting[caption={\texttt{StreamingSystem}}]{apendices/prompts/StreamingSystem.txt}

El último mensaje del usuario se sustituye por la siguiente plantilla, donde \texttt{\{0\}} son los fragmentos recuperados (título, enlace público, organismo, fecha y texto de cada uno) y \texttt{\{1\}} es la pregunta:

\lstinputlisting[caption={\texttt{StreamingUserTemplate}}]{apendices/prompts/StreamingUserTemplate.txt}

La plantilla pide descartar solo los documentos que no tratan el tema de la pregunta y aprovechar los demás. La versión usada en las ejecuciones del estudio de fragmentación (capítulo~\ref{cap:estudio-chunking}) era distinta: pedía ignorar todo el contexto si no coincidía exactamente con el tema y responder a partir del historial. Se cambió tras el estudio, porque se sospechaba que causaba las abstenciones de las planas, pero al repetir las ejecuciones con la versión actual las abstenciones apenas cambiaron (sección~\ref{sec:comprobaciones}). La lista de fuentes se añade siempre que se han recuperado documentos, se usen o no (sección~\ref{sec:impl-rag}).

\subsection{Respuesta en \emph{streaming} sin documentos}

Cuando la búsqueda no devuelve documentos (por ejemplo, en un saludo o en una pregunta de seguimiento), la conversación continúa con este \emph{prompt}, que prohíbe dar de memoria datos concretos de una norma:

\lstinputlisting[caption={\texttt{ConversationalSystem}}]{apendices/prompts/ConversationalSystem.txt}

\section{Pipeline jerárquica}

\subsection{Expansión de la consulta}

Antes de buscar, la pregunta se reformula en terminología administrativa y se separa en una consulta léxica y otra semántica. \texttt{\{0\}} es la consulta del usuario y \texttt{\{1\}} la conversación anterior.

\lstinputlisting[caption={\texttt{QueryExpansion}}]{apendices/prompts/QueryExpansion.txt}

\subsection{Generación en lenguaje claro}

La respuesta se genera con el modelo de la plataforma a partir del texto completo de los documentos padre recuperados (\texttt{\{0\}}). La regla~6 y la última línea de \texttt{StreamingSystem} se añadieron durante la puesta a punto del estudio de fragmentación, al comprobar que las respuestas incluían fórmulas de cortesía y consejos generales que no se pueden comprobar en los documentos y que bajaban la fidelidad medida por RAGAS (sección~\ref{sec:validez}).

\lstinputlisting[caption={\texttt{ClearLanguageGeneration}}]{apendices/prompts/ClearLanguageGeneration.txt}

\section{Ingesta}

Durante la ingesta jerárquica, cada fragmento hijo (\texttt{\{0\}}) se enriquece con dos preguntas que un ciudadano haría y que el fragmento responde, antes de generar su \emph{embedding}.

\lstinputlisting[caption={\texttt{FragmentEnrichment}}]{apendices/prompts/FragmentEnrichment.txt}
